% =====================================================================
\section{Conclusion}
\label{sec:conclusion}
% =====================================================================

We have given a complete and self-contained treatment of the chain
\[
  \text{non-commutative matrix rings}
  \;\longrightarrow\;
  \text{cyclic codes}
  \;\longrightarrow\;
  \text{Galois duality and hulls}
  \;\longrightarrow\;
  \text{Gray maps}
  \;\longrightarrow\;
  \text{quantum codes}.
\]

The algebraic core is that over $R=M_2(A_t)$ the Galois twist must be the
anti-automorphism $\theta_s=\sigma_s\circ(\cdot)^{\mathsf T}$, so that the dual
of a right ideal is a right ideal with the correct dimension
(Theorems~\ref{thm:duality}, \ref{thm:dualformula}, \ref{lem:failure}). The hull
theory is developed for arbitrary --- in particular non-principal --- right
ideals, and the natural conjecture that Galois hulls are principal is false
(Theorem~\ref{thm:nonprincipal}), with the phenomenon confined to $t\ge2$ and
$p\mid n$.

The coefficient expansion transports the Galois pairing to the form
$u^{[s]}(I_{4n}\otimes H_t)v^{\mathsf T}$ (Theorem~\ref{thm:gram}); duality and
hulls are transported exactly, non-principal hulls included; and the image of a
Galois hull is self-orthogonal \emph{without} any reflexivity hypothesis
(Theorem~\ref{thm:hullso}), which is what makes the whole chain work in the
general Galois regime. A structural obstruction rules out the Euclidean
coordinate change in characteristic two with $t$ even
(Theorem~\ref{thm:obstruction}), and its hypotheses are exact.

On the quantum side, three routes (R1--R3) are available, their parameters and
distances are given with the hypotheses visible, and the distance is
characterised as the minimum weight of a quotient code
(Theorem~\ref{thm:cosetdistance}), yielding the coset Singleton bound
$K+d\le\Nlen+1$ (Theorem~\ref{thm:cosetsingleton}). A feasibility inequality
(Theorem~\ref{thm:feasibility}) describes the entanglement budget, and the
measured frontier is reported (Proposition~\ref{prop:frontier},
Table~\ref{tab:budget}).

The best codes obtained are $[[24,8,4]]_9$ among stabilizer codes and
$[[20,2,10;18]]_4$ and $[[24,2,10;18]]_9$ among entanglement-assisted codes.
\textbf{No record parameter is claimed.} The stabilizer codes are one or two
below the best known distances of \cite{GrasslCodetables}, and the
entanglement-assisted codes cannot be compared at all, since no
entanglement-assisted table exists. What is claimed is a new \emph{origin} for
quantum codes: non-commutative Galois hulls, transported exactly, with
non-principal hulls admitted and with two new bounds on the resulting
parameters.

\paragraph{Open problems.}
(i) Determine the Galois hull \emph{dimension} distribution in the wild
families, where the cyclic-code enumeration has resisted complete treatment for
$15$ families.
(ii) Find a closed or recursive description of the set of weights of the
quotient code of Theorem~\ref{thm:cosetdistance} in terms of the CRT data of
$C$.
(iii) Prove or refute the two conjectures stated at the end of
Section~\ref{sec:hulls}: $\Hull_s(C)\subseteq\sigma_{2s}(C)$, and
$\dim\Hull_s(C)=\dim\Hull_s(C^{\perp_s})$.
(vi) Decide the sharper form {\rm (O-3.1b)} of the refuted conjecture
(Remark~\ref{rem:o31b}): every witness currently known is a code that is
itself non-principal, so the question whether a \emph{principal} cyclic code
can have a non-principal Galois hull in the wild regime is open.
(iv) Determine whether the obstruction of Theorem~\ref{thm:obstruction} is the
only one: we have not proved that a congruence exists for all $t$ odd in
characteristic two, only that it exists for $t=1,3,5$ at $q=4,8$.
(v) Extend the comparison to $q=8$ and to entanglement-assisted codes, neither
of which is tabulated.
